% Source: 04 -Mon Oct 5 2026.pdf, page 2. Content remains in source order.
\sourcepage{04}{2}
\section{Boolean Circuits}
In 1971 Cook proved the first NP-completeness result on a problem in boolean logic. We'll work with an equivalent variant defined on boolean circuits.

We only look at combinatorial circuits, so acyclic interconnection of boolean gates, because a cycle in a circuit implies that the circuit has memory. 

\subsection{Boolean gates} A gate is an elementary switching circuit of three types:
\begin{enumerate}
\item AND
\sourcefigure[0.25]{assets/04/page-002-and-gate.png}
output wire $y=1\Longleftrightarrow$ input wires $x_1,x_2,\ldots,x_k$ all carry 1. Thus $y=\bigwedge_{i=1}^k x_i$.
\item OR
\sourcefigure[0.25]{assets/04/page-002-or-gate.png}
output wire $y=1\Longleftrightarrow$ at least one input wire carries 1. Thus $y=\bigvee_{i=1}^k x_i$.
\item NOT
\sourcefigure[0.175]{assets/04/page-002-not-gate.png}
output wire $y=1\Longleftrightarrow$ input wire $x=0$. Thus $y=\neg(x)$ ($=\bar x$).
\end{enumerate}

Some subsets create a Elementary universal family of gates 

\subsection{Boolean Circuit}
\textbf{DEF:} A Boolean Circuit (BC) is an acyclic interconnection of gates, where the output of a gate can become input to one or more other gates.

Acyclic because we are interested in combinatorial circuits 